\section{Introduction}
\label{sec:intro}

Aircraft and ships are tracked by cooperative surveillance systems that
rebroadcast positions computed by an on-board Global Navigation Satellite
System (GNSS) receiver. In aviation the carrier is Automatic Dependent
Surveillance--Broadcast (ADS-B) \citep{faa2007adsb,khan2025adsbsurvey}; at sea
it is the Automatic Identification System (AIS) \citep{itu2014ais}. Neither
channel authenticates the position it carries, so a receiver that is jammed
or spoofed passes the corrupted position directly to air traffic control
(ATC) and vessel traffic services (VTS). Interference has become a routine
operational condition rather than a rare event. Aviation monitoring reports a
sustained rise in GNSS anomalies in and around conflict zones
\citep{aireon2025,felux2025navigating,ieropoulos2024impact,lo2025observations}.
Maritime trackers reported roughly 700 jamming and spoofing events per day
across the global fleet in 2024 \citep{windward2025}. In the Baltic Sea, the
number of affected vessels grew from about 1,225 in the first quarter of 2025
to more than 5,800 in the second \citep{gpspatron2025}. A single two-week
window in mid-2025 recorded more than 12,000 spoofing incidents involving
over 3,000 vessels \citep{kpler2025}, and sustained spoofing has been
documented in the Persian Gulf \citep{skuld2025,maritimeexec2025}.

Detecting interference is necessary but not sufficient. Once a detector
flags an aircraft or a vessel, its satellite-derived positions must be
discarded, yet the operator still has to deconflict traffic, enforce
separation and verify that the asset stays clear of restricted areas. In a
narrow ship channel or a busy terminal area, losing a reliable track for
even half an hour is operationally significant. The engineering problem is
therefore \emph{post-detection reconstruction}: estimating where the asset is
from its last trusted state, from signals that do not depend on satellites,
and from what similar assets have done before.

Today this gap is usually filled by trajectory predictors used as
extrapolators. These predictors are almost always domain-specific. Aircraft
predictors are trained on flights \citep{pang2019lstm,kwak2025evolution}, and
vessel predictors are trained on AIS voyages
\citep{li2023ship,shin2024aisacnet,nguyen2024transformer}. Each must be
rebuilt, retuned and revalidated when it is moved to a new region, a new
traffic mix or a new transport mode. For an organisation that operates both
airspace and coastal surveillance, or for a vendor that supplies both, this
means maintaining two unrelated model families for what is essentially the
same task.

In the conference version of this work \citep{felendler2026trajrag} we
introduced TrajRAG for aircraft. TrajRAG converts trajectory segments into
textual state descriptions and stores them in phase-specific vector indices.
During a blackout it retrieves phase-consistent historical segments and lets
a large language model (LLM) propose the next position, which is accepted
only inside a kinematic feasibility envelope. That study was limited to
aviation: it did not establish whether the design is specific to aircraft
dynamics and ADS-B, or whether it is a reusable reconstruction method. This
paper answers that question. Our thesis is that \emph{nothing in the
reconstruction engine needs to be learned per domain}: moving to a new
transport mode should require only a new prompt template, a new retrieval
index and a small set of physically motivated constants, all derived from
domain data by fixed procedures.

We test this thesis on shipping, which differs from aviation in almost
every respect that matters to a predictor. Vessels are an order of
magnitude slower. AIS is sampled irregularly. Phase dynamics are dominated
by berthing manoeuvres rather than altitude changes, the environmental
drivers are currents and waves rather than winds aloft, and the
GNSS-independent signal available on board is different. If the same
engine, with its LLM frozen and no maritime gradient update, competes with
dedicated maritime models, the result is evidence that retrieval and
prompting can replace retraining in this class of engineering problem.

\paragraph{Contributions} The contributions of this paper are:
\begin{enumerate}
\item \textbf{A domain-agnostic formulation (AI contribution).} We separate
post-detection reconstruction into a shared engine (retrieval, frozen-LLM
proposal, feasibility guard and feedback loop) and a domain adapter of eight
elements. We specify the adapter formally and show that it is the only part
that changes between aviation and shipping (Section~\ref{sec:framework}).
\item \textbf{A maritime instantiation built without training.} We
construct the maritime adapter (AIS state schema, gyrocompass heading as the
trusted signal, departure/transit/arrival phases, marine environmental
context and a voyage prompt template), together with a new public-data
corpus of 3,506 voyages from 26.2 million AIS reports in two United States
port corridors (Sections~\ref{sec:instantiation} and
\ref{sec:data}).
\item \textbf{A rigorous maritime benchmark.} We compare against fourteen
baselines: classical, recurrent, retrieval-only, seven modern forecasters and
AIS-ACNet, a trained ship-specific network published in this journal
\citep{shin2024aisacnet}. All are run on a common held-out test set with
paired significance tests and distance- and time-normalised errors
(Section~\ref{sec:results}).
\item \textbf{A transferable calibration procedure for the feasibility
guard.} We show that the procedure used to calibrate the guard transfers
across domains, while the constant it produces does not. Vessel dynamics
give a phase-dependent displacement-ratio distribution with no counterpart
in aviation (Section~\ref{sec:calibration}).
\item \textbf{A real-time dual-domain system design.} We describe how one
deployment serves ATC and VTS, and we show that the reconstruction latency
fits inside the reporting interval of both surveillance channels
(Sections~\ref{sec:architecture} and \ref{sec:latency}).
\end{enumerate}

\paragraph{Relation to the conference version} The conference paper
\citep{felendler2026trajrag} covers aviation only. The domain-adapter
formulation, the maritime instantiation and corpus, the full maritime
evaluation (fourteen baselines, normalised errors, significance tests and
corridor analysis), the cross-domain calibration analysis and the
dual-domain system architecture are new. Aviation results are summarised in
Section~\ref{sec:aviation-reference} only as the reference point for the
cross-domain comparison; we do not repeat the aviation ablations, backbone
study or real-world case studies, and refer the reader to the conference
paper for them.

The rest of the paper is organised as follows. Section~\ref{sec:related}
reviews related work. Section~\ref{sec:framework} presents the
domain-agnostic framework and the system architecture. Section~\ref{sec:instantiation}
describes the two instantiations. Section~\ref{sec:design} details the
experimental design, and Section~\ref{sec:results} reports the results.
Section~\ref{sec:discussion} discusses implications and threats to validity,
and Section~\ref{sec:conclusion} concludes.
